%% training.tex
%% Chapter: Efficient Training of Multimodal Large Language Models.
%% Designed to be \input{} into a master ACM acmart manuscript.
%% Bibliography file: training.bib

\section{Efficient Training of Multimodal Large Language Models}
\label{sec:efficient-training}

\subsection{Motivation: The Cost of Multimodal Training}
\label{subsec:training-motivation}

Training a multimodal large language model is expensive at every stage. A standard LLaVA-style pipeline~\cite{Liu202304828,Liu202310999} already requires days of training on dozens of A100 GPUs; scaling to the billion-image corpora and 70B-parameter backbones used in frontier systems such as InternVL2~\cite{Chen202413948} demands weeks of compute accessible only to a few organizations. The bottleneck compounds at three levels: \emph{pretraining}, where the model must align visual and linguistic representations across massive corpora; \emph{instruction fine-tuning}, where the model must follow human intent across diverse tasks; and \emph{downstream adaptation}, where a pretrained backbone must transfer to new applications with minimal overhead. A substantial body of research therefore focuses on reducing the cost of each stage without sacrificing quality.

This section organizes that research into four complementary directions. Section~\ref{subsec:efficient-pretrain} reviews efficiency-oriented pretraining strategies that reduce the data and compute needed to build capable visual representations. Section~\ref{subsec:efficient-sft} covers efficient instruction fine-tuning methods that minimize the parameters updated and the annotations consumed. Section~\ref{subsec:rethink-stages} examines principled reassessments of how training stages should be structured, motivated by the observation that the conventional sequential pipeline inherits its design from text-only LLMs without multimodal justification. Section~\ref{subsec:peft} surveys parameter-efficient transfer learning techniques that enable task-specific adaptation with a fraction of the full model's parameters.

\subsection{Efficiency-Oriented Pretraining}
\label{subsec:efficient-pretrain}

The dominant MLLM pretraining recipe freezes a CLIP encoder, trains a lightweight connector on image--caption pairs, and optionally unlocks more parameters for instruction tuning. This pipeline is empirically convenient but rests on choices---which encoder, which connector, how much data, whether to freeze the language model---that are rarely questioned. Three works challenge these defaults and, by doing so, substantially reduce the compute required to reach a given quality level.

Idefics2~\cite{Laurençon202402246} approaches pretraining as a controlled ablation study. Starting from an open reproduction of Flamingo~\cite{Alayrac202204228}, the authors systematically vary the visual backbone (CLIP versus SigLIP), the connector architecture (cross-attention versus perceiver resampler versus linear projection), the training data mixture (pairs only versus interleaved documents), and whether the language model is frozen or fully fine-tuned. Two findings drive most of the quality gain: \emph{unfreezing the language model during pretraining} and using a \emph{perceiver resampler} with SigLIP. The perceiver resampler compresses patch tokens to a fixed-length sequence, which reduces the context fed to the LLM and makes unfreezing feasible without memory explosion---an example of connector design and LLM-freezing choice interacting rather than being independent. The resulting model, trained entirely on open data (OBELICS~\cite{Laurençon202309462} and LAION-5B subsets), matches proprietary systems of comparable size across standard benchmarks. Crucially, Idefics2 also shows that the interleaved-versus-pairs data mix matters less than the architectural and freezing decisions, which challenges the common practice of investing heavily in interleaved data collection as a first step.

VILA~\cite{Lin202312533} reaches the opposite conclusion on data distribution, but the two findings are not in conflict. VILA investigates which pretraining data enables \emph{in-context learning} (ICL)---the ability to condition on multiple images at inference time---rather than merely improving single-turn accuracy. Three controlled experiments show that interleaved image--text documents are necessary for ICL, while image--caption pairs alone are not, because interleaved text teaches the model to retrieve visual context from earlier positions in the sequence. VILA further shows that unfreezing the LLM during interleaved pretraining is essential: a frozen LLM cannot learn to integrate multi-image context regardless of data format. The efficiency gain is compositional---a balanced mix of CC3M, SBU, COCO, and a curated subset of MMC4~\cite{Zhu202304341} achieves competitive ICL performance with far fewer tokens than OpenFlamingo, because distribution rather than volume drives the capability.

TinyLLaVA~\cite{Zhou202402289} shifts focus entirely to the small-scale regime (1B--4B parameters) and introduces a framework that decouples three orthogonal axes: visual encoder (CLIP-L, SigLIP-So400M, DINOv2), language model (Phi-2, Gemma-2B, TinyLlama), and training recipe (data mixture, resolution, learning rate). By holding two axes constant while sweeping the third, the authors show that data quality dominates model scale: a 3.1B TinyLLaVA model trained on carefully curated ShareGPT4V~\cite{Chen202311972} captions matches or exceeds LLaVA-1.5-7B on six standard benchmarks, using less than half the parameters. The framework also reveals that the visual encoder matters more than the language model at this scale, because SigLIP-So400M provides a more stable training signal than CLIP-L under dense-caption supervision. TinyLLaVA thus establishes a practical design principle for resource-constrained pretraining: invest in a stronger visual encoder and higher-quality captions rather than a larger language model.

Several concurrent works reinforce and extend these insights. MobileVLM~\cite{Chu202312539} and MobileVLM V2~\cite{Chu202401130} replace the LLM backbone with MobileLLaMA and introduce a lightweight downsampling projector that reduces visual token count fourfold, enabling deployment on mobile devices without sacrificing MLLM-style instruction following. LLaVA-Phi~\cite{Zhu202401101} substitutes Phi-2 for Vicuna and demonstrates that a 2.7B model with high-quality instruction data matches 7B counterparts on several benchmarks, echoing TinyLLaVA's data-quality thesis at an even smaller scale. Bunny~\cite{He202402116} further shows that the language backbone's pretraining corpus affects downstream VLM quality more than the connector architecture, suggesting that the community's attention to connector design may be disproportionate relative to the payoff.

Across these works, two recurring conclusions stand out. First, the decision to unfreeze or freeze the language model during pretraining has an outsized effect on alignment quality, and the right choice depends on the connector: compressive connectors (perceiver resampler, Q-Former) enable unfreezing; token-preserving connectors (linear projection) do not. Second, data distribution and quality govern capability acquisition more than raw token volume, which means that curating a smaller but richer pretraining corpus is a more efficient path to a capable model than scaling a na\"ively collected one.

\subsection{Efficient Instruction Fine-Tuning}
\label{subsec:efficient-sft}

Instruction fine-tuning transforms a pretrained MLLM into a model that follows natural-language directives. Na\"ively fine-tuning all parameters risks overwriting the language capabilities acquired during pretraining and is computationally expensive. Two strategies dominate the efficient fine-tuning literature: adapter-based methods that inject small trainable modules into a frozen backbone, and hypernetwork-based methods that generate adapter weights dynamically from each input.

The adapter approach was established in the text-only domain by Houlsby adapters~\cite{Houlsby201902110} and popularized in the LLM era by LLaMA-Adapter~\cite{Zhang202302301}, which injects zero-initialized attention adapters at the upper transformer layers, concentrating adaptation where the language model is most sensitive to instruction format. LLaMA-Adapter V2~\cite{Gao202303747} extends this design to vision by adding a visual projection layer and joint visual--language instruction data, demonstrating that top-layer adapters are sufficient for both modalities with minimal interference to lower-layer representations. LaVIN~\cite{Luo202315023} takes the adapter idea further by framing visual instruction fine-tuning as a modality-routing problem. It inserts Mixture-of-Modality Adaptation (MMA) modules---lightweight bottleneck adapters with a soft gating mechanism that amplifies visual information when processing image tokens and suppresses it for text tokens---into a frozen LLaMA backbone. The gating design lets a single set of adapter weights specialize for both modalities without dedicated modality-specific branches, and the resulting system trains in approximately 1.4 hours on eight A100 GPUs while matching LLaVA-7B on ScienceQA.

LaVIN's MMA modules use fixed adapter weights regardless of the input image, which limits their effectiveness when visual content varies substantially across inputs---scenes, documents, charts, and fine-grained objects each impose different alignment demands on the LLM. HyperLLaVA~\cite{Wen202403447} addresses this by replacing fixed adapters with a \emph{dynamic visual expert}: a small hypernetwork that takes patch features from the visual encoder and generates layer-specific adapter weight deltas conditioned on the current image. Because the hypernetwork output is a low-rank delta, the total parameter overhead remains modest, but the generated weights now specialize to each input. HyperLLaVA outperforms LaVIN on benchmarks involving charts, documents, and scene understanding, where the required alignment between visual features and LLM representations shifts across inputs. The trade-off is inference latency: generating adapter weights adds a forward pass through the hypernetwork, approximately doubling visual encoding time while leaving LLM decoding unchanged.

Beyond MMA and hypernetwork adapters, instruction-aware visual selection provides a complementary efficiency strategy. InstructBLIP~\cite{Dai202305519} extends the frozen BLIP-2 Q-Former by feeding the instruction as an additional input, causing the Q-Former to select the subset of patch tokens most relevant to the current instruction before passing them to the LLM. This filtering step reduces the effective visual context length dynamically and improves task performance precisely on queries that require focused spatial attention. Otter~\cite{Li202306943} pursues yet another angle, showing that in-context instruction tuning---providing a few demonstrations at fine-tuning time rather than single-turn examples---substantially improves instruction generalization, especially for multi-step visual reasoning tasks.

The collective picture from these methods is that \emph{selective} adaptation consistently outperforms na\"ive full fine-tuning in the low-resource regime, and the choice of selection mechanism---gating (LaVIN), dynamic generation (HyperLLaVA), query-based filtering (InstructBLIP), or in-context examples (Otter)---should match the nature of the task distribution. Fixed adapters suffice when the visual domain is narrow; dynamic adapters pay off when visual content is heterogeneous; query-based filtering is most effective when tasks require precise spatial reasoning; and in-context tuning is most valuable when the instruction distribution is broad.

\subsection{Rethinking Training Stages for Efficiency}
\label{subsec:rethink-stages}

The conventional MLLM training pipeline applies three sequential stages: visual--language alignment pretraining, supervised instruction fine-tuning, and optional preference alignment. This structure is largely inherited from text-only LLMs without systematic justification for the multimodal setting. Three works challenge the conventional stage design from different angles: SPHINX-X~\cite{Gao202402935} argues for task unification through joint training across stages, Cobra~\cite{Zhao202403520} eliminates the transformer backbone's quadratic attention bottleneck entirely through architectural replacement, and TinyGPT-V~\cite{Yuan202316862} demonstrates that a carefully sequenced minimal design can reach usable quality on a single GPU within one day.

SPHINX-X extends the SPHINX model~\cite{Lin202310250} into a family spanning 400M to 13B parameters and trains jointly on a corpus that spans visual question answering, optical character recognition, chart understanding, document analysis, grounding, and mathematical reasoning simultaneously, rather than across separate stages. The curriculum assigns sampling weights proportional to each task's volume and adjusts them dynamically based on a held-out validation mixture, which detects and corrects task forgetting on the fly. Two efficiency gains follow from joint training. First, the model encounters each capability at every training step, preventing the catastrophic forgetting that sequential stages induce. Second, shared visual representations transfer across tasks: grounding annotations improve spatial attention, which benefits VQA tasks that require precise region identification. Additionally, SPHINX-X introduces a token compression scheme that merges adjacent visual tokens via learned pooling, reducing sequence length by 4$\times$ without measurable accuracy loss on most benchmarks and directly addressing the quadratic attention cost at high resolution.

Cobra~\cite{Zhao202403520} targets a more fundamental efficiency bottleneck: the transformer's self-attention mechanism scales quadratically with sequence length, so processing many visual tokens is inherently costly regardless of how well the training stages are organized. Cobra replaces the LLM backbone with Mamba~\cite{Gu202312110}, a selective state-space model with $O(n)$ inference complexity, and integrates dual visual encoders---DINOv2~\cite{Oquab202309896} for semantic features and SigLIP for language-aligned features---through a lightweight MLP projector. The resulting system retains the standard two-stage training procedure (alignment pretraining on LLaVA-595K, then instruction fine-tuning on LLaVA-665K) but benefits from recurrent rather than attention-based decoding. Cobra-3B matches LLaVA-1.5-7B on standard benchmarks while reducing inference latency by approximately 3$\times$ at 1024 tokens and 5$\times$ at 2048 tokens. The efficiency gain grows with sequence length, making Cobra particularly attractive for high-resolution or video inputs where token counts are large. Unlike adapter-based methods, Cobra's efficiency is architectural and therefore not reversible: the gain is unconditional, but compatibility with the broader transformer ecosystem is sacrificed.

TinyGPT-V~\cite{Yuan202316862} explores the lower bound of the MLLM design space: a 2.7B model built on Phi-2~\cite{Li202312048} and trained from scratch on a single A100 GPU in under 24 hours using a four-stage progressive unfreezing schedule. Stage one warms up the connector on image--caption alignment; stage two pretrain the connector and visual encoder jointly on a curated mix of LLaVA and MiniGPT-4~\cite{Zhu202302506} data; stage three applies instruction fine-tuning on visual QA; stage four adds grounding through multi-task fine-tuning. At each boundary, the previously frozen backbone components are selectively unlocked, preventing interference between the stages. Two design choices underpin the efficiency: a Q-Former~\cite{Li202301596} that compresses 256 patch tokens to 32 query tokens before passing them to the LLM, and Phi-2's code-heavy pretraining, which provides structured reasoning without requiring many instruction-tuning examples. TinyGPT-V approaches the benchmark scores of 7B models trained for a week, demonstrating that the minimum viable MLLM is smaller and cheaper than the literature commonly assumes.

The three works identify orthogonal levers for stage-level efficiency: joint multi-task optimization (SPHINX-X), architectural backbone replacement (Cobra), and progressive staged unfreezing (TinyGPT-V). Combining these levers---for instance, training a Mamba-based model with joint multi-task objectives and a progressive unfreezing schedule---is a natural but largely unexplored direction for future work.

\subsection{Parameter-Efficient Transfer Learning}
\label{subsec:peft}

Parameter-efficient transfer learning (PEFT) methods adapt a pretrained model to a new task by updating only a small subset of parameters. Originally developed for text-only LLMs, PEFT methods require non-trivial extension to multimodal settings because the visual encoder, connector, and LLM each contribute differently to task performance, and the visual encoder introduces modality-specific alignment demands that text-only methods ignore.

Low-Rank Adaptation (LoRA)~\cite{Hu202106865} is the most widely adopted PEFT baseline. It inserts trainable rank-$r$ matrix pairs ($\mathbf{W} = \mathbf{W}_0 + \mathbf{A}\mathbf{B}$, where $r \ll \min(d, k)$) into the attention projections of each transformer layer and updates only those pairs during fine-tuning. Applied to the LLM backbone of an MLLM with the visual encoder frozen, LoRA works well when the target task is close to the instruction-tuning distribution but degrades on visually demanding tasks---fine-grained classification, document understanding, spatial grounding---because the frozen visual encoder's features do not adapt to the new domain. Visual Prompt Tuning (VPT)~\cite{Jia202201986} addresses this complementary gap by prepending a small number of learned prompt tokens to the visual encoder's patch sequence, steering the encoder's attention toward task-relevant features without modifying its weights. VPT-Deep, which injects prompts at every transformer layer rather than only the first, outperforms VPT-Shallow substantially on distribution-shifted inputs, suggesting that multi-level visual steering is necessary for large domain gaps. AdaLoRA~\cite{Zhang202203985} further refines LoRA by dynamically reallocating the rank budget across layers using singular value decomposition, pruning ranks in layers that contribute little gradient signal and redirecting them to high-impact layers; for multimodal tasks, this consistently outperforms fixed-rank LoRA at the same total parameter count because visual understanding concentrates representation changes in a small subset of attention heads. QLoRA~\cite{Dettmers202305441} reduces the memory footprint of LoRA fine-tuning by quantizing the frozen backbone to 4-bit NormalFloat, enabling full PEFT runs for 13B models on a single 48GB GPU and democratizing MLLM adaptation for practitioners without large-scale infrastructure.

Two methods specifically designed for multimodal PEFT address limitations that generic LoRA and VPT cannot resolve. EAS~\cite{Zong202318083} introduces a modality-aware adapter design: bottleneck adapter layers are inserted into the LLM's feed-forward sublayers with a gating mechanism that amplifies the adapter's contribution when processing visual tokens and suppresses it for textual tokens, allowing one set of adapter weights to specialize asymmetrically across modalities. EAS further includes a \emph{Selector} component that allocates adapter capacity based on gradient signal: a short 100-step pilot run identifies which layers carry the largest gradient magnitude for the target task, and only those layers receive full adapter modules while the remaining layers receive only the gating mechanism without the bottleneck projection. On average, EAS finds that the top 40\% of LLM layers dominate the gradient signal for visual tasks, consistent with the LLaMA-Adapter finding that upper layers are the primary site of instruction adaptation. The combination of modality-aware gating and gradient-based layer selection allows EAS to match full fine-tuning on VQAv2 and OKVQA while updating fewer than 2\% of backbone parameters.

MEMVP~\cite{Zeng202409752} targets the memory bottleneck that arises when applying visual prompt tuning to high-resolution inputs. Standard VPT prepends $p$ prompt tokens to the patch sequence, growing the sequence length by $p$ regardless of input resolution; at high resolutions---documents at 1024$\times$1024 or medical images at 2048$\times$2048---the patch count already strains GPU memory, and additional prompt tokens make fine-tuning infeasible. MEMVP resolves this by moving prompt information from the sequence dimension to the channel dimension: instead of prepending tokens, it adds a learned bias to the patch embeddings at each layer, parameterized as the outer product of a spatial map and a channel vector. The full bias tensor is $\mathcal{O}(h \times w \times d)$, matching the resolution of the patch grid, but because it is expressed as a rank-1 outer product, the trainable parameters are only $\mathcal{O}(h + w + d)$---orders of magnitude fewer than standard VPT. MEMVP matches or exceeds VPT-Deep on six benchmarks while reducing fine-tuning memory by 8$\times$, and, because channel-dimension prompts do not lengthen the sequence, it composes naturally with LoRA on the LLM side: visual encoder adaptation via MEMVP and LLM adaptation via LoRA can be trained jointly with a combined parameter count below 1\% of the backbone.

Multimodal-specific extensions continue to broaden the PEFT landscape. VL-Adapter~\cite{Sung202112227} inserts shared adapter modules into both the visual encoder and the LLM and shows that sharing adapter weights across modalities improves generalization to held-out tasks, because the shared parameters encode cross-modal alignment rather than modality-specific representations. SoRA~\cite{Ding202312023} replaces a single low-rank matrix pair with a sparse mixture of small expert pairs selected per token, capturing the heterogeneity of visual features more faithfully than a shared delta. Together, EAS, MEMVP, VL-Adapter, and SoRA establish that multimodal PEFT is not simply text-only PEFT applied to a larger pipeline: the visual encoder, connector, and LLM each have distinct adaptation characteristics, and methods that respect these differences through modality-aware gating, memory-efficient channel prompting, or adaptive rank allocation consistently outperform uniform approaches.

\subsection{Cross-Cutting Comparison and Open Challenges}
\label{subsec:training-summary}

The four lines of work reviewed in this chapter address complementary stages of the MLLM training pipeline and together provide a toolkit for building capable multimodal models within realistic resource budgets. Table~\ref{tab:training-compare} summarizes representative methods across the four dimensions.

\begin{table}[t]
  \caption{Representative efficient training methods across four dimensions. ``Params updated'' reports the fraction of backbone parameters modified. ``Key lever'' identifies the primary source of efficiency gain.}
  \label{tab:training-compare}
  \small
  \setlength{\tabcolsep}{4pt}
  \begin{tabular}{@{}p{0.17\linewidth}p{0.20\linewidth}p{0.14\linewidth}p{0.36\linewidth}@{}}
    \toprule
    Method & Stage & Params updated & Key lever \\
    \midrule
    Idefics2~\cite{Laurençon202402246}   & Pretraining      & 100\%   & Ablation-driven connector + LLM unfreezing \\
    TinyLLaVA~\cite{Zhou202402289}       & Pretraining      & 100\%   & Data quality over model scale \\
    VILA~\cite{Lin202312533}             & Pretraining      & 100\%   & Interleaved data for in-context learning \\
    LaVIN~\cite{Luo202315023}            & Instruction FT   & $<$4\%  & Modality-gated MMA adapters \\
    HyperLLaVA~\cite{Wen202403447}       & Instruction FT   & $<$2\%  & Per-image dynamic adapter generation \\
    SPHINX-X~\cite{Gao202402935}         & Stage design     & 100\%   & Joint multi-task training with token compression \\
    Cobra~\cite{Zhao202403520}           & Stage design     & 100\%   & $O(n)$ Mamba backbone \\
    TinyGPT-V~\cite{Yuan202316862}       & Stage design     & 100\%   & Progressive staged unfreezing \\
    EAS~\cite{Zong202318083}             & Adaptation       & $<$2\%  & Modality-aware gating + gradient-based selection \\
    MEMVP~\cite{Zeng202409752}           & Adaptation       & $<$1\%  & Channel-dimension visual prompting \\
    \bottomrule
  \end{tabular}
\end{table}

Three open challenges cut across all four directions. First, no comprehensive study has jointly optimized across all four stages: a model might be pretrained with VILA's interleaved recipe, instruction-tuned with LaVIN adapters, structured as a SPHINX-X joint curriculum, and domain-adapted with EAS, but whether these choices compose or conflict is unknown. The interaction between pretraining data distribution and downstream PEFT effectiveness is particularly unclear: EAS and MEMVP assume a well-pretrained backbone, but when the backbone is pretrained on a narrow distribution (as in TinyGPT-V's single-GPU regime), the gradient signal that EAS relies on for layer selection may be unreliable.

Second, architectural efficiency and parameter efficiency address different bottlenecks and are rarely studied together. Cobra demonstrates that replacing the attention mechanism with Mamba yields inference gains that no amount of adapter pruning can match; however, Mamba's recurrent structure is less amenable to VPT-style prompt injection than transformer attention, and combining Cobra-style backbones with EAS or MEMVP-style adaptation remains unexplored.

Third, scaling laws for efficient MLLMs are largely absent. The literature documents that efficiency techniques preserve quality at a given scale but does not characterize how the efficiency gap changes as models grow from 3B to 70B parameters. Whether small-model lessons---data quality over scale (TinyLLaVA), joint multi-task training (SPHINX-X)---transfer to larger regimes, or whether they are artifacts of the low-resource setting, is an open empirical question that systematic scaling experiments could resolve.

%% End of \section{Efficient Training of Multimodal Large Language Models}
